% Part: counterfactuals
% Chapter: introduction
% Section: paradoxes-material

\documentclass[../../../include/open-logic-section]{subfiles}

\begin{document}

\olfileid{cnt}{int}{par}

\olsection{في مفارقات الشرط المادي}

كان سي.~آي. لويس في طليعة من أنكروا أن يكون~$!A \lif !B$ تمثيلًا
صحيحًا لكل قول إنجليزي من نحو «إذا \dots، فإن \dots». وخصّ بالنقد
استعمال الشرط المادي في كتاب وايتهد وراسل \emph{مبادئ الرياضيات}، إذ
كانا يقرآن~$\lif$ «يستلزم». واعتراضه في موضعه إن حُمّل~$\lif$ هذا
المعنى: فكل قضية كاذبة~$p$ تستلزم عندئذ أن~$p$ تستلزم~$q$، لأن
$p \lif (p \lif q)$ تصدق إذا كذبت~$p$؛ وكل قضية صادقة~$q$ تستلزم
كذلك أن~$p$ تستلزم~$q$، لأن~$q \lif (p \lif q)$ تصدق إذا صدقت~$q$.

والمعلوم عند المناطقة أن \emph{الاستلزام} المنطقي ليس رابطًا، بل علاقة
بين !!{formula} أو العبارات. فلا ينبغي إذن أن نقرأ
$\lif$ بمعنى «يستلزم»، اتقاءً للالتباس.\footnote{ما زالت قراءة
  «$\lif$» بمعنى «يستلزم» واسعة الانتشار بين الرياضيين وعلماء الحاسوب،
  مع أن الفلاسفة يتقون اللبس الذي كشفه لويس بقراءته «فقط إذا».} فإذا
تحاشينا تلك القراءة زال أصل اعتراض لويس: فليست~$p$ «مستلزمة» لـ~$q$
لمجرد أن الرمز~$p$ ناب عن جملة إنجليزية كاذبة في الواقع. والحكم في
$p \Entails q$ مردّه إلى \emph{جميع} !!{valuation}؛ ولذلك يبقى
$p \Entails/ q$ وإن صادف أن الجملة الممثلة بـ~$p$ كاذبة.

على أن الغرابة تبقى في بعض عبارات «إذا \dots، فإن \dots»، ومنها مثال
لويس:
\begin{quote}
إذا كان القمر مصنوعًا من الجبن الأخضر، فإن~$2+2=4$.
\end{quote}
وكذلك في الاستدلالين:
\begin{quote}
  القمر ليس مصنوعًا من الجبن الأخضر. ولذلك، إذا كان القمر مصنوعًا
  من الجبن الأخضر، فإن~$2+2=4$.

  $2+2 = 4$. ولذلك، إذا كان القمر مصنوعًا
  من الجبن الأخضر، فإن~$2+2=4$.
\end{quote}
ومع ذلك، إن لم تكن «إذا \dots، فإن \dots» إلا~$\lif$، صدقت الجملة
ولزم صحة الاستدلالين من غير قيد.

ومثال آخر ناشئ من القضية الصادقة منطقيًا
$(!A \lif !B) \lor (!B \lif !A)$. فإن أخذنا اعتباطًا جملتين خبريتين
$S$ و$T$ من صحيفة، قضت هذه القضية بصدق القول: «إذا كانت~$S$ فإن~$T$،
أو إذا كانت~$T$ فإن~$S$».

\end{document}



